\section{Introduction}\label{sec:intro}

Strain-gradient elasticity adds to the classical energy a quadratic form in the strain
gradient, governed in Mindlin's form~II by a sixth-order tensor
$\tD$~\cite{Mindlin1964,MindlinEshel1968}. For crystals it is no longer unknown: Admal et al.~\cite{Admal2017} derived its atomistic
representation, and Lazar et al.~\cite{Lazar2022} and Po et al.~\cite{Po2019} published
complete cubic tensors of fcc and bcc metals from interatomic potentials, which now enter non-singular theories of dislocations~\cite{Po2018,Lazar2024};
granular micromechanics gives anisotropic gradient tensors in terms of the
fabric~\cite{Pirmoradi2024}. Using such tensors for heterogeneous bodies
requires the operations that classical micromechanics performs on
fourth-order tensors: products, inverses, orientation averages and concentration tensors.
Inverses of sixth-order tensors appear in the inclusion and inhomogeneity problems of
Monchiet and Bonnet~\cite{Monchiet2011a,Monchiet2013,Monchiet2015} and of Ma et al.~\cite{Ma2018},
the compliance in estimates of Reuss type~\cite{Nazarenko2021a}, and all four operations
in the self-consistent scheme of Solyaev~\cite{Solyaev2023}.

For fourth-order tensors, Walpole~\cite{Walpole1984} gave bases with multiplication tables
for the 32 crystal classes, which reduce these operations to products of small matrices;
bases of the same kind also give the orientation averages of third-order elastic
constants~\cite{Kube2026}, and harmonic bases of tensors of any order, with tensorial
texture coefficients, follow from the representations of the rotation
group~\cite{Krause2025}. For gradient tensors, the corresponding algebra is known for
isotropy~\cite{Monchiet2011a,Auffray2015}, also without the minor
symmetries~\cite{Monchiet2011b}, and for transverse isotropy~\cite{Monchiet2013}, where
it was built to solve inhomogeneity and crack problems~\cite{Monchiet2015}. The anisotropic tensors have been classified. Olive and Auffray~\cite{Olive2013}
obtained the 17 symmetry classes of the sixth-order tensor, and Auffray et
al.~\cite{Auffray2013} wrote the matrix of each class in an orthonormal basis, and later
for the 48 classes of the whole law, fifth-order coupling included~\cite{Auffray2019}. Le Quang
et al.~\cite{LeQuang2021} classified the tensors by their symmetry planes, Auffray et
al.~\cite{Auffray2021} made the harmonic structure explicit in two dimensions, and
Poorsolhjouy~\cite[table~B.2]{Poorsolhjouy2026} counted, class by class, the tensors
without major symmetry. Demonteaux et al.~\cite{Demonteaux2025} report an orthogonal basis for the cubic
tensor, obtained from the character and projection formulas of the cubic group; its
formulas are not yet published. For the cubic class, Lazar et al.~\cite{Lazar2022} wrote
the tensor with eleven constants, reduced its matrix to diagonal blocks of sizes five and
three, and obtained its eigenvalues by factoring the characteristic polynomial of the
block of size five into a quadratic and a cubic, conditions of positive definiteness and Voigt
average. Averages of the same Voigt type have been computed for fcc metals in Mindlin's
second-gradient theory~\cite{Bagherpour2024a,Bagherpour2024b}. Nazarenko et
al.~\cite{Nazarenko2021a} gave the inverse of the coupled constitutive law in tensorial
form, explicitly for hemitropic materials.

These works establish how many constants each class has and where they sit in a matrix:
they are parametrisations. An algebra is a different object, a basis in which the
product of two tensors of the class is a product of small blocks, so that the inverse,
the orientation average and the concentration tensor follow block by block. To our
knowledge no such basis, and no closed-form compliance or Reuss estimate, has been
published for the cubic or the hexagonal class. For the cubic class an orthogonal basis
has been announced~\cite{Demonteaux2025}; what we add there are its explicit formulas,
the multiplication table and their uses. The hexagonal class has a further interest.
In two dimensions, the gradient tensor of a hexagonal lattice is anisotropic whereas its
classical tensor is isotropic, which makes wave propagation anisotropic at high
frequency~\cite{Rosi2016}. In three dimensions the hexagonal classical tensor is
transversely isotropic, and the single gradient modulus that breaks transverse isotropy
has been identified only as the coefficient of a fixed matrix~\cite{Auffray2013}, which
we write $\eta_{\mathrm A}$ (their $\eta$) to keep $\sgrad$ for the strain gradient.

This paper constructs these algebras and uses them. First, we give coordinate-free bases,
with closed multiplication tables, of the algebras of the classes $\clscub$ and $\clshex$
on the strain-gradient space, of dimensions 15 and 32. Their symmetric elements are the
gradient tensors, with 11 and 22 constants, and the other elements are the
non-reciprocal tensors of the same classes. Repeated irreducible components are labelled
by the chain $\GL\supset\SOrth\supset G$, with the conversion to the labels of
Lazar~\cite{Lazar2016}. The hexagonal modulus absent from transverse isotropy is the
splitting of the part of azimuthal order $\pm3$ of the totally symmetric harmonic
component into two irreducible parts, equivalently, up to a factor, the component of
$\tD$ along the hexagonal structure tensor $\mathrm{Re}\,(\bm a+\iu\,\bm n\times\bm a)^{\ot6}$, and it
equals $4\eta_{\mathrm A}$. The hexagonal algebra is the transversely isotropic algebra of
Monchiet and Bonnet~\cite{Monchiet2013} with one of its idempotents split in two; what is
new there is the orthonormal basis, the identification of the extra modulus and the closed
forms. Second, the gradient compliance $\tD^{-1}$ follows in closed block form,
with conditions of invertibility and positive definiteness block by block and the change of basis to and from the constants of Lazar et
al.~\cite{Lazar2022} and the matrices of Auffray et al.~\cite{Auffray2013}. Third, the
Voigt, Reuss and Hill estimates of untextured polycrystals are weighted means of the
block moduli. The Reuss estimate lies below the Voigt one, with equality only for an
isotropic crystal; in the cubic case their vector-like blocks differ by a matrix of rank
one, and the hexagonal extra modulus drops out of the Voigt estimate. The cubic formulas
are applied to the published atomistic data of fcc and bcc metals. Finally, as an illustration, the self-consistent
scheme of Solyaev~\cite{Solyaev2023} with the complete cubic tensor reduces to the
inversion of four blocks and, for grains of physical size, sees the grain only through
its Voigt estimate. On the way we correct the average over the sphere of the sixth-order
Eshelby tensor of Ma et al.~\cite{Ma2018}, on which that scheme relies.

The averages are algebraic estimates, not bounds, and the paper is not a homogenisation
theory: the heterogeneity of the classical stiffness alone induces gradient moduli that
grow with the square of the grain size~\cite{Malyi2021}, and no average of $\tD$ contains
them. By Malyi's estimate they equal the Voigt estimate of $\tD$ at a grain diameter of
about $\NumAlInducedCrossoverNm$~nm for Al and $\NumCuMEAMInducedCrossoverNm$~nm for Cu
and dominate above it (section~\ref{sec:averages-cubic}). The averages are therefore the
part of the effective tensor carried by $\tD$: an isotropic substitute for the crystal
tensor, as in~\cite[section~7]{Lazar2022}, the gradient tensor of nanocrystals whose grains
are a few lattice parameters across, and an ingredient of a scheme that also accounts for
the heterogeneity of $\tC$. No hexagonal gradient tensor of a crystal has been published, so the hexagonal
results are illustrated with the granular micromechanics of Pirmoradi
et al.~\cite{Pirmoradi2024}, in which the extra modulus is the sixfold harmonic of the
contact distribution.

Section~\ref{sec:setting} recalls form~II and builds the isotropic algebra as a template,
sections~\ref{sec:cubic}--\ref{sec:application} treat the cubic and hexagonal algebras,
the compliance with the data and the verification, the averages and the illustration,
and the appendices convert our labels to those of Lazar and derive the averaged Eshelby
tensor.
